% !TeX root = ../papers/nc-p3.tex
\section{横断する因子対の自由積表示と正則性}
\label{sec:transverse-factor-amalgam}
\begin{theorem}[横断する二平面での完全判定]\label{thm:transverse-factor-amalgam}
定理\ref{thm:transverse-product-even-rank}の四評価階数3・横断性の下で、
標準Koszul AS正則となる必要十分条件は$\kappa_U\ne0$である。
これは$U$の可換単位積の非結合性と同値で、$V$で検査しても同じである。
この条件は二つの反対鉛筆の被約性を問わない。
\end{theorem}
\begin{proof}
$\kappa_U=0$なら前節の指標の関係が純テンソルになる。
以後両方の$\kappa$が非零とする。
一つの積の三つの対称関係$\lambda(p*q)$ ($\lambda\in U^*$)が
階数1の形式を含まないことを示す。もし含めば単位$e$での縮約は
非零の$\lambda$だから、規格化して$\lambda(p*q)=\lambda(p)\lambda(q)$となる。
これは全ての結合子を消す。$L_p$を左乗法とすると
$\langle [L_p,L_q]r,s\rangle$は前節の面積形式の積に符号を掛けたものである。
$\bar p,\bar q$が独立ならこの交換子の階数は2、核は$ke$である。
$[L_p,L_q]r$も指標で消えるから$\lambda$の計量双対は$ke$に入る。
ところが$\langle e,p*q\rangle$は非退化な計量そのもので矛盾する。
従って三つの対称関係を課す三生成代数$B_U$は三変数二次完全交叉の
Koszul双対で、AS三次元となる。$B_V$も同じである。

$B_U$の三文字のうち二文字空間$W$を単位を消す線型形式全体とする。
三関係の線型結合で$W\otimes W$に入るものは単位評価で消えるので零である。
従って全対称積の像は$\operatorname{Sym}^2 W$の像と同じ三次元である。
これらの像は中心的である。実際、$\operatorname{Sym}^2 W$の基底を
独立な三つの次数2の中心変数に置き、全ての反交換子をこれらの対応する線型結合とする
Clifford代数へ$B_U$は全射する。Cliffordの順序付き平方自由基底から次数級数は
$(1+t)^3/(1-t^2)^3=(1-t)^{-3}$となり、$B_U$の級数と等しいため同型である。

二文字の部分代数$D$の三次関係は各語の最初と最後の文字の交換である。
実際、二つの平方は中心的で、二文字の反交換子も中心的である。
この交換だけで、各語の偶位置・奇位置を別々に整列できる。長さ$2m,2m+1$の次元は
それぞれ$(m+1)^2,(m+1)(m+2)$となる（交換関係は各軌道の語を同一視するだけである）。
Clifford基底で見た二文字部分も三つの中心変数上階数4で同じ級数なので、
上の三次関係で尽くされる。$B_U$は$D$上左右に自由で基底は$1,Z_U$である。
残り一文字$Z_U$をClifford基底で先頭または末尾に置けばよい。

第一・第二位置での二文字空間はそれぞれ$X^*,Y^*$である。
$U$表示と$V$表示の二文字座標は違っても、三次の両端を交換する関係空間は
両端に同じ、中央に独立な基底変換で保たれる。従って両側の$D$は同じ添字付き代数である。
六関係は$B_U,B_V$の各三関係の合併だから全代数は$A=B_U *_D B_V$となる。
各位置で$D$の右係数を持つ、$Z_U,Z_V$が交互に現れる語が基底である。
これを確かめるには、同じ側の二文字を掛けてその側の自由基底で書き、
間の$D$係数を右へ送る。重なりは同じ側の三文字、または同じ側の二文字と係数に限り、
その側の結合性と自由基底の一意性で解消する。反対向きでも同じ議論が通る。
従って$A$は両$B$上でも左右に自由である。

各頂点の自明加群を一緒に$k_*$と書く。誘導加群には短完全列
\[
 0\longrightarrow A\otimes_D k_*\longrightarrow
 (A\otimes_{B_U} k_*)\oplus(A\otimes_{B_V} k_*)\longrightarrow k_*\longrightarrow0
\]
がある。最初の写像は二つの商写像の差である。非空交互語はその末尾と異なる側の商でのみ
基底として残り、空語では対角の差が増大の核を張る。
$D$の自明加群の最小解消の階数は$1,2,2,1$、内部次数は$0,1,3,4$である。
実際、長さ$n$の語は二つの位置類の文字の別々の対称冪である。
先頭一文字の乗法の核はその位置類の交代二文字と残る対称冪で表され、
三次の二関係を掛けて全て得られる。その乗法の核は第二の位置類でも交代二文字を
取ったもので、一つの四次関係の自由倍となる。これは二次元空間での乗法
$W\otimes\operatorname{Sym}^j W\to\operatorname{Sym}^{j+1}W$の核が
$\bigwedge^2W\otimes\operatorname{Sym}^{j-1}W$であることを二回使ったものである。

自由性から誘導加群に対する$\operatorname{Tor}^A(k_*,-)$は
それぞれ$D,B_U,B_V$でのTorとなる。長完全列で$D$から両$B$への写像は
次数0と1では対角に単射、次数2と3では内部次数の相違で零である。
従って$A$の自明加群のTorは各位置で対角次数に集中し、階数は$1,4,6,4,1$で尽くされる。
元のGramテンソルは全ての一文字切断で階数4である。
実際、$\beta(x,-)=0$なら各単位での評価により$a(x)=m(x)=0$となり、
両者の核は別の直線である。右側も同様である。
隣接階数6と上のKoszul性から定理\ref{thm:frobenius-tensor-reconstruction}の
三つの条件を満たす。双対の完全対合と左右Koszul解消から標準AS正則性を得る。
\end{proof}

\begin{corollary}[印のある横断チャートの同型]
正則な二表示の間で順序付きの二因子を保つ添字保持同型は、
両文字空間の線型同型と直交二因子の相似変換の組で、
二つの表示$(a(x)*_U b(f),m(x)*_V n(f))$を対応させるものの存在と同値である。
二つの相似比は共通とする。印を忘れるときは標的の全ての横断する順序付き因子対を試す。
\end{corollary}
\begin{proof}
定理\ref{thm:unique-linear-factor-partner}で二位置ずつの変換が揃い、
全射Gram表示の同値から共通相似比の変換を得る。逆は関係列への代入で従う。
\end{proof}

\subsection{根基が直線の場合の共通な重み付き多項式環}
\begin{theorem}[根基直線の完全判定]\label{thm:radical-factor-amalgam}
補題\ref{lem:product-radical-cubic}の四評価階数3と$U\cap V=kr$を仮定する。
標準Koszul AS正則となる必要十分条件は、二文字関係$t$が非退化であることである。
\end{theorem}
\begin{proof}
六次元Gram空間の各元との対合を関係とみなす。
$t$は非零の二行二列関係であり、$U,V$の各三次元関係は
$S_U,S_V$の極形式そのものである。軸評価は各三文字空間に全射だから、
これらの階数1の関係は全窓にも純テンソルを与える。これで必要性を得る。
逆に$t$が非退化ならどちらの極形式系にも階数1の形式はない。
実際、$S_U(-,e,r)=0$なので階数1の対称形式があれば、その一次因子は$e$か$r$を消す。
$S_U(l,e,-)=0$なら$l\in kr$で、$S_U(r,-,-)$は階数2である。
$S_U(l,r,-)=0$なら$t$の非退化性から$l\in ke$で、$S_U(e,-,-)$も階数2である。
$e,r$が比例する場合も通る。$V$側も同様である。
三関係を各々に課した三生成代数を$B_U,B_V$とする。
全射Gram表示の双対の単射性により三関係は各々独立である。
前定理と同じ二次完全交叉の双対性により、どちらも標準のKoszul AS三次元である。

まず一方の三文字を単位$e$に関して分け、$W=\operatorname{ann}(e)$を二文字とする。
$S_U(l,e,-)$は$l$の退化計量による像なので、
二文字だけによる関係は$l\in kr$に対応するただ一式$t$である。
従って$\operatorname{Sym}^2 W$の像は二次元であり、全対称積の像は三次元である。
前定理と同じClifford表示を使う。すなわち全対称積の像の基底を
三つの独立な中心変数としてClifford代数を作ると、全射とHilbert級数の一致で同型である。
最初の二つの中心変数は二文字から取れる。このため二文字部分は非退化な一関係$t$だけで
表示される代数$D$になる。実際、対応する二文字のClifford基底は二中心変数上階数4で、
$1/(1-t_0)^2$の級数を与える（$t_0$は級数変数）。
非退化な二文字の一関係代数も各位置で基底を取り直せば二文字多項式環の級数となる。
第三の中心二次式$Q_U$を任意に補えば、$B_U$は$D[Q_U]$上左右に
$1,Z_U$を基底として自由である。これは二文字を先に、または後に並べたClifford基底による。
$V$側でも同じである。$D$の表示は第一・第二位置で$t,t^{\mathsf t}$であり、実際に同じ部分代数である。

ここで二側の第三の中心式の取り方が重要になる。$\langle s,r\rangle=1$なる
$s\in\mathcal H$を選ぶ。$U$上の対称双線型形式$q_U$を
$q_U(e,-)=\langle s,-\rangle|_U$として延長する。$q_U(e,r)=1$だが、
全ての極形式と二文字の二次式は$(e,r)$で消えるので、$q_U$の類を
第三の中心変数$Q_U$に取れる。$V$上でも同じように$q_V,Q_V$を選ぶ。
$q_U,q_V$を$(a,b),(m,n)$で引き戻すと、六つ目の関係は
\[
 \langle s,\beta(x,f)\rangle=q_U(a(x),b(f))+q_V(m(x),n(f))+b_0(\bar x,\bar f).
\]
実際、いずれかが軸のとき第一・第二項はそれぞれ指定した値になり、差は二文字だけの式となる。
反転位置では全てが転置される。

$D$については両位置を二次元空間に同定して$t$を非零交代形式とすることができる。
次位置の転置は符号だけ違うので、この同定で$D$は添字付き二変数多項式環である。
$b_0$と次位置での転置はこの可換環で同じ二次式を与える。
$Q_U$とその転置は三文字の対称な中心式を偶奇の表示で書いたもので、
二文字に対して$Q_{U,p}w_{p+2}=w_p Q_{U,p+1}$と交換する。
$V$でも同じである。従って$D$に次数2の中心変数$Q$を加えた共通部分
$D'=D[Q]$を、一方で$Q_U$、他方で$-Q_V-b_0$へ送って埋め込める。
後者は変数の平行移動で$D[Q_V]$と同じ部分代数となる。
位置を揃えた$D'$は重み$(1,1,2)$の多項式環である。
六関係のうち五つが両三関係の合併（$t$を共有）で、残りは
この$Q$の二表示の同一視であるから、$A=B_U*_{D'}B_V$である。

前定理の交互語基底と誘導加群の短完全列を$D'$に適用する。
$D'$のTorは内部次数$1,1,2$の三元上の外積代数である。
$\operatorname{Tor}_1$で二つの一次元方向は両側へ対角に単射、二次の生成元は零へ行く。
$\operatorname{Tor}_2$の内部次数2の一元は二文字の定義関係$t$に対応し、
各$B$の独立な二次関係の一つへ送られるため単射である。
これは二つの一次文字を含めたbar複体の二文字サイクルをそのまま写せば分かる。
内部次数3の二元と$\operatorname{Tor}_3$の内部次数4の一元は内部次数の差で零へ行く。
従って長完全列は再び$A$の対角Torだけを残し、その階数は$1,4,6,4,1$である。
各一文字切断の単射性は$a,m$と$b,n$の核直線の相違から従い、
Gramテンソルの完全対合により標準AS正則性を得る。
\end{proof}

\begin{corollary}[両対称反転窓の有限な分類手順]\label{cor:bisymmetric-effective-recognition}
隣接階数6の両対称テンソルの標準AS性と、正則なものの添字保持同型を有限な代数的条件で判定できる。
\end{corollary}
\begin{proof}
反対階数6では完全対称化可能性と対称完全交叉の条件を検査する。
階数6と8以外は定理\ref{thm:bisymmetric-six-or-eight}で排除される。
階数8では両鉛筆が四次完全交叉であることを確認する。
両側が被約なら系\ref{cor:all-reduced-apolar-recognition}のグラフ判定を使う。
その同型候補は射影的にも有限型の方程式で扱える。具体的には抽象グラフを
$\mathbb P^3\times\mathbb P^3$のHilbertスキームで探し、
両射影の同型性、二つの次数配分の一致および切断の制限条件を課す。
グラフの$(1,1)$に関するHilbert多項式はどれも$8n$である。

一側が非被約ならその鉛筆には階数2以下の元がある。平方の場合は
軸評価を検査し定理\ref{thm:double-plane-ore}を使う。
相異なる因子なら両側の四軸の評価階数3を検査し、
軸像の交わりが二次元以上のものは補題\ref{lem:reciprocal-apolar-factors}の
二行二列の純関係で排除する。残る二場合は本節の二定理で決まる。
同型は定理\ref{thm:unique-linear-factor-partner}から両文字空間での二つの基底変換と
六次元Gram空間の相似変換による係数同値と同値である。
必要なら印を指定した標的の因子を全て動かす。従って各条件は有限変数の方程式・不等式、
または指定された固定多項式のHilbertスキームで決まる。
この判定の対象は両対称表示を持つ反転窓である。
\end{proof}
